\section{高考考频分析}

\subsection{近五年考点分布（2021--2025 新课标Ⅰ卷）}

\begin{center}
\begin{tabular}{c|c|c|c|c|c}
\textbf{知识模块} & \textbf{2021} & \textbf{2022} & \textbf{2023} & \textbf{2024} & \textbf{2025} \\ \hline
物质的量 & 3 & 3 & 3 & 3 & 3 \\
氧化还原 & 3 & 3 & 0 & 0 & 3 \\
离子反应 & 3 & 3 & 3 & 3 & 3 \\
周期律与周期表 & 6 & 6 & 6 & 6 & 6 \\
反应热与盖斯定律 & 3 & 3 & 3 & 3 & 3 \\
化学反应速率与平衡 & 6 & 3 & 8 & 6 & 6 \\
电解质溶液 & 8 & 11 & 9 & 6 & 11 \\
电化学 & 3 & 3 & 3 & 6 & 3 \\
非金属元素 & 3 & 6 & 6 & 6 & 6 \\
金属元素 & 3 & 3 & 3 & 3 & 3 \\
有机化学 & 15 & 15 & 15 & 15 & 15 \\
化学实验 & 10 & 10 & 13 & 13 & 10 \\
物质结构（选考） & 15 & 15 & 15 & 15 & 15 \\
工业流程 & 13 & 13 & 10 & 10 & 13 \\
反应原理综合 & 13 & 13 & 13 & 13 & 10 \\ \hline
\textbf{总计} & 100 & 100 & 100 & 100 & 100
\end{tabular}
\end{center}

\textit{注：数字为分值（含选考）。必考部分总分 85 分，选考 15 分，合计 100 分。}

\subsection{必考板块（每年必出）}

\[
\begin{array}{c|c|c}
\text{模块} & \text{年均分值} & \text{考查形式} \\ \hline
\text{元素周期律与元素周期表} & 6 分 & 选择题 1--2 题 \\
\text{电解质溶液} & 8--11 分 & 选择题 + 原理题结合 \\
\text{有机化学} & 15 分 & 选考题（或选择题 1--2 题） \\
\text{化学实验} & 10--13 分 & 实验探究大题 \\
\text{工业流程} & 10--13 分 & 工业流程大题 \\
\text{化学反应原理} & 10--13 分 & 原理综合大题 \\
\text{阿伏加德罗常数} & 3 分 & 选择题 \\
\text{离子方程式} & 3 分 & 选择题 \\
\end{array}
\]

\subsection{常考板块（高频轮考）}

\[
\begin{array}{c|c}
\text{模块} & \text{考查特点} \\ \hline
\text{电化学} & 近 3 年出现频率高，多出现选择题 \\
\text{化学反应速率与平衡} & 常结合图像出现在原理题中 \\
\text{元素化合物} & 分散在工业流程和实验题中考查 \\
\text{氧化还原反应} & 每年都有，但独立选择题出现频率降低 \\
\end{array}
\]

\subsection{轮考板块}

\[
\begin{array}{c|c}
\text{模块} & \text{说明} \\ \hline
\text{热化学} & 有时独立出选择题，有时融入原理题 \\
\text{化学与STSE} & 偶有选择题，分值 3 分 \\
\text{化学计量} & 计算能力要求逐年提高 \\
\end{array}
\]

\subsection{命题趋势与复习策略}

\textbf{趋势分析}：
\begin{enumerate}
  \item \textbf{情境化}：试题情境越来越丰富（科研成果、工业生产、环境保护）
  \item \textbf{综合化}：同一题目跨模块考查（如元素化学 + 原理 + 计算）
  \item \textbf{实验能力}：实验题分值稳定在 10--13 分，重视设计能力和误差分析
  \item \textbf{计算能力}：盖斯定律、$K$ 相关计算、产率计算要求准确
  \item \textbf{阅读量}：题干变长，信息提取能力要求提高
\end{enumerate}

\textbf{复习优先级}：
\[
\begin{array}{c|c|c}
\text{优先级} & \text{模块} & \text{建议投入} \\ \hline
\star\star\star & \text{实验探究题} & 20\% \\
\star\star\star & \text{工业流程题} & 15\% \\
\star\star\star & \text{化学反应原理} & 15\% \\
\star\star\star & \text{选考题（有机/结构）} & 15\% \\
\star\star & \text{选择题专题（阿伏伽德罗/离子反应/周期律）} & 15\% \\
\star\star & \text{元素化合物} & 10\% \\
\star & \text{化学与STSE、基础概念} & 10\% \\
\end{array}
\]
